\documentclass[12pt, a4paper]{article}
\usepackage[utf8]{inputenc}
\usepackage{hyperref}
\usepackage{listings}
\usepackage{xcolor}
\usepackage{geometry}

\geometry{margin=1in}

\definecolor{codegreen}{rgb}{0,0.6,0}
\definecolor{codegray}{rgb}{0.5,0.5,0.5}
\definecolor{codepurple}{rgb}{0.58,0,0.82}
\definecolor{backcolour}{rgb}{0.95,0.95,0.92}

\lstdefinestyle{htmlcode}{
    backgroundcolor=\color{backcolour},   
    commentstyle=\color{codegreen},
    keywordstyle=\color{magenta},
    numberstyle=\tiny\color{codegray},
    stringstyle=\color{codepurple},
    basicstyle=\ttfamily\footnotesize,
    breakatwhitespace=false,         
    breaklines=true,                 
    captionpos=b,                    
    keepspaces=true,                 
    numbers=left,                    
    numbersep=5pt,                  
    showspaces=false,                
    showstringspaces=false,
    showtabs=false,                  
    tabsize=2
}

\lstset{style=htmlcode}

\title{Integration Guide: Embedding Skills Portal into Main Domain via Iframe}
\author{Technical Implementation Document}
\date{\today}

\begin{document}

\maketitle

\section{Overview}
This document outlines the steps required to integrate the Skills portal website, hosted separately at \url{https://ksaw.vercel.app}, into the main domain \url{https://vtu.ac.in} under the path \url{https://vtu.ac.in/ksaw}.

To avoid merging the two codebases or modifying complex server reverse proxy routing rules, we use a clean, full-screen \texttt{iframe} layout. This ensures that the Skills portal is loaded within the main website structure while keeping the domain address in the browser address bar unchanged.

\section{Implementation Steps}

\subsection{Step 1: Create the Host Page on the Main Domain}
On the main website codebase (\url{https://vtu.ac.in}), create a new directory or page handler at the path corresponding to \texttt{/ksaw}.

\subsection{Step 2: Add the Embed Code}
Depending on the technology stack of the main website, embed the code snippet shown below.

\subsubsection{Option A: For Standard HTML/Vanilla PHP Websites}
Save the following code as \texttt{index.html} (or equivalent layout file) inside the \texttt{/ksaw} directory on the main website server:

\begin{lstlisting}[language=HTML, caption=Full Screen HTML Iframe Wrapper]
<!DOCTYPE html>
<html lang="en">
<head>
    <meta charset="UTF-8">
    <meta name="viewport" content="width=device-width, initial-scale=1.0">
    <title>Skills Form</title>
    <style>
        body, html {
            margin: 0;
            padding: 0;
            height: 100%;
            overflow: hidden;
        }
        iframe {
            width: 100%;
            height: 100%;
            border: none;
        }
    </style>
</head>
<body>
    <iframe src="https://ksaw.vercel.app" title="Skills Form"></iframe>
</body>
</html>
\end{lstlisting}

\subsubsection{Option B: For React/Next.js/TSX Websites}
If the main website is a modern React application, use the following React Component:

\begin{lstlisting}[language=HTML, caption=React / TSX Iframe Component]
import React from 'react';

export default function KsawPage() {
  return (
    <div style={{ width: '100vw', height: '100vh', overflow: 'hidden', margin: 0, padding: 0 }}>
      <iframe
        src="https://ksaw.vercel.app"
        title="Skills Form"
        style={{ width: '100%', height: '100%', border: 'none' }}
      />
    </div>
  );
}
\end{lstlisting}

\section{Verification}
Once deployed, visit \url{https://vtu.ac.in/ksaw}. Confirm that:
\begin{enumerate}
    \item The browser address bar remains exactly as \url{https://vtu.ac.in/ksaw}.
    \item The page correctly renders the Skills application form hosted on Vercel.
    \item The page fits the entire window size without secondary scrollbars.
\end{enumerate}

\end{document}
